\chapter{设计总结}\label{ch:limit}

\section{总体架构}

\code{AFDX_top} 采用单一 125~MHz 时钟，以发送核心、接收核心、服务通道、配置模块和时间基准组成端系统。发送和接收两侧都按每拍一字节组织数据通路，逐 VL 配置与状态存放在块 RAM 中，状态机按固定阶段完成读出、判决和写回。全设计通过 19 个 \code{afdx_ram_sdp} 实例承载载荷、配置、状态和协议文件，模块层次见表~\ref{tab:hierarchy}。

发送核心把消息存储与线路发送分开。分页缓存接纳应用消息，发送引擎按 VL 的 BAG 和队列状态调度，组帧器预读载荷并生成协议字段，双网 MAC 共用状态机和发送时刻。提前候选和帧头预备使允许时刻到来前完成控制准备；线路争用及其配置条件由第~\ref{sec:budget}~节的预算给出。

接收核心为 A、B 两网分别设置 MAC、解析器和环形缓冲，共享 VL 判决与应用交付。载荷先试写，经过帧格式、配置、完整性和冗余判决后提交，拒绝时回退。每网独立保存序号状态，冗余管理决定有效副本的交付，异常分类与统计事件记录各帧的处理结果。

\section{服务与主机控制}

RTC、SNMP 和 615A 通过服务通道共享收发核心，服务报文与应用报文按各自 VL 配置参与发送调度和接收检查。RTC 在报文握手处采样时间戳，分拍完成时间格式换算与偏差计算；SNMP 通过请求与回应 RAM 完成 BER 解析、对象访问和回应组装；615A 以 TFTP 传输协议文件及加载区数据，并通过状态文件报告进度与结果。各项协议的处理范围和字段检查见第~\ref{ch:svc}~章。

主机通过配置口管理 VL、服务参数和统计计数器。配置写在引擎许可的周期写入，配置读与计数器取样有固定时序；带宽、排队、缓存和交付约束由主机配置共同满足，编程步骤见第~\ref{ch:cfg}~章。顶层 GMII 字节流、事件端口和状态寄存器提供运行观察接口，其定义见第~\ref{ch:if}、\ref{ch:verif}~章。

按上述结构，端系统以 2\,442\,112 位逻辑存储和逐拍一级组合深度的数据通路达到设计指标：每网每方向载荷容量 956.44~Mbps；线路空闲时端到端时延上界 26.360~µs、发送抖动为 0；争用时抖动由式~(\ref{eq:jitter}) 给出，满足配置约束时不超过 50~µs；128 条发送 VL 与 4096 条接收 VL 可全部配置并同时启用。综合资源、时序检查、功能测试和性能测量结果集中列于测试报告。
